%----------------------------------------------------------------------------------------
%	SECTION TITLE
%----------------------------------------------------------------------------------------

\cvsection{Experiencia}

%----------------------------------------------------------------------------------------
%	SECTION CONTENT
%----------------------------------------------------------------------------------------

\begin{cventries}

%------------------------------------------------

\cventry
{Técnico de Diagnóstico - Área de ingeniería de pruebas} % Job title
{Jabil Circuit} % Organization
{Zapopan, Jalisco} % Location
{Julio 2018 - Febrero 2020} % Date(s)
{ % Description(s) of tasks/responsibilities
\begin{cvitems}
\item {Lectura e interpretación de "log files" porporcionados por el equipo de prueba ICT.}
\item {Conocimiento y uso de instrumentos de medición como multímetros digitales, osciloscopios digitales, medidores LCR, generadores de señales y máquina de rayos X}.
\item {Experiencia en el área de ICT realizando pruebas de depuración a las tarjetas a través del ICT, así como el diagnóstico a las tarjetas.}
\item {Experiencia en el área de FVT, inspeccionando y probando ensambles o subensambles provenientes de las líneas de producción, así como su diagnóstico a nivel ensamble o componente. }
\item {Desarrollo de reportes de análisis de fallas, recolectando y analizando la información para identificar la causa raíz del problema y enviarle el reporte al cliente.}
\item {Habilidad para leer y entender diagramas esquemáticos, hojas de datos de componentes, lista de materiales (BOMs), procedimientos técnicos y dibujos de diseño de PCB.}
\end{cvitems}
}

%------------------------------------------------    

\cventry
{Programador Analista (PA) - Área de IT} % Job title
{Jabil Circuit} % Organization
{Zapopan, Jalisco} % Location
{Febrero 2020 - Marzo 2022} % Date(s)
{ % Description(s) of tasks/responsibilities
\begin{cvitems}
\item {Desarrollo y mantenimiento de aplicaciones en C\# con .NET framework y Visual Basic para diferentes clientes dentro de la compañía.}
\item {Corrección de errores y bugs detectados en aplicaciones ejecutadas en el ambiente de producción (PRD).}
\item {Experiencia y desarrollo en el área de Front-End utilizando lenguajes como JavaScript, así como HTML5, CSS utilizando el framework de Angular.}
\item {Experiencia y desarrollo en el área de Back-End utilizando lenguajes como C\#.}
\item{Uso del marco de trabajo de Scrum para la realización de los proyectos.}
\item {Desarrollo de manuales de usuario y materiales de entrenamiento, así como manuales y diagramas técnicos.}
\item {Uso de DevOps para automatzar el trabajo de desarrollo de software.}
\item {Experiencia con bases de datos relaciones con SQL Server y PostgreSQL.}
\item {Diseño de diagramas de entidad-relación para el diseño y modelado de bases de bases de datos relacionales.}
\item {Conocimiento del ciclo de vida del desarrollo del software (SDLC).}
\item {Familiaridad con plataformas en la nube como AWS.}
\item {Cración y manejo de web services con XML y REST (JSON, XML).}
\item {Creación y manejo de APIs para la comunicación entre diferentes aplicaciones o sistemas.}
\end{cvitems}
}

%------------------------------------------------

\cventry
{Business System Analyst (BSA) - Área de IT} % Job title
{Jabil Circuit} % Organization
{Zapopan, Jalisco} % Location
{Marzo 2022 - Noviembre 2024} % Date(s)
{ % Description(s) of tasks/responsibilities
\begin{cvitems}
\item {Estar en constante comunicación con el Product Owner (PO) para formar los requerimientos del producto final (la aplicación)}.
\item {Constante comunicación con el cliente o stakeholders a tavés de reuniones presenciales o en línea para la toma de requerimientos.}
\item {Uso de devOps para registrar y documentar las historias de usuario y seguimiento del proyecto.}
\item {Realización de documentación como el BRD y FSD en el cual se trasladaba la información del cliente en requerimientos técnicos para el desarrollador.}
\item {Identificación y creación de requerimientos funcionales, no funcionales y criterios de aceptación de la aplicación, con base en los requerimientos del cliente.}
\item {Diseño de mokups utilizando Figma.}
\item {Generación de diagramas del proceso.}
\item {Generación de diagramas UML para la planificación de información.}
\item {Desarrollo de test cases diseñados para validar los requerimientos funcionales y no funcionales del sistema.}
\item{Desarrollo de pruebas funcionales para verificar y validar que la aplicación funcione de acuerdo a los requerimientos.}
\item {Manejo de web services y APIs para la comunicación entre diferentes aplicaciones o sistemas.}
\item {Diseño de diagramas de entidad-relación para el diseño y modelado de bases de bases de datos relacionales.}
\item {Conocimiento de Jira para la gestión de procesos.}
\item {Uso de bases de datos para la consulta de información a través de diferentes tablas.}
\item {Implementación de la aplicación en el ambiente de producción.}
\item{Uso del marco de trabajo de Scrum para la ralización de proyectos a través de sprints.}
\end{cvitems} 
}

%------------------------------------------------


%------------------------------------------------

\end{cventries}